% ============================================================
\section{理性偏好的结构性质：拟线性}\label{sec:quasilinear}
% ============================================================

拟线性是应用工作（产业组织、公共经济学、福利分析）中使用最频繁的效用形式。
本节给出它的\textbf{偏好侧完整刻画}：三个序条件 $\iff$ 线性计价物表示，
并推导其两大后果（无财富效应、福利的精确货币度量）。

\subsection{设定与三个计价物条件}

设 $X=\R\times Y$：第一坐标 $t\in\R$ 为\textbf{计价物}（数量可正可负、连续可分），
$Y$ 为一般集合（不必是欧氏空间），$\succeq$ 为 $X$ 上的弱序。

\begin{condition}[计价物有价值]\label{cond:num-valuable}
对一切 $y\in Y$ 与 $t,t'\in\R$：
\[
(t,y)\succeq(t',y)\ \Longleftrightarrow\ t\ge t'.
\]
同一 $y$ 下，计价物越多越好，且严格。
\end{condition}

\begin{condition}[补偿可行]\label{cond:compensable}
对任意 $y,y'\in Y$，存在 $t\in\R$ 使
\[
(0,y)\ \sim\ (t,y').
\]
任意两个 $Y$ 方案之间的价值差都能用计价物\textbf{恰好补偿}。
\end{condition}

\begin{condition}[无财富效应]\label{cond:no-wealth}
若 $(t,y)\succeq(t',y')$，则对一切 $d\in\R$：
\[
(t+d,\,y)\ \succeq\ (t'+d,\,y').
\]
给两个方案同时增加等量计价物，不改变比较结果——偏好对计价物的“初始存量”不敏感。
\end{condition}

\begin{example}[无财富效应的直觉]
若 $(10,\,y)\sim(0,\,y')$——“持有 10 单位计价物配 $y$”与“不持计价物配 $y'$”无差异，
则同时加 100：$(110,\,y)\sim(100,\,y')$ 仍无差异。
财富水平的变化不翻转、不破坏任何比较。
\end{example}

\subsection{拟线性表示定理}

\begin{theorem}[拟线性表示定理]\label{thm:ql}
设 $\succeq$ 为 $X=\R\times Y$ 上的弱序，满足
条件~\ref{cond:num-valuable}、\ref{cond:compensable}、\ref{cond:no-wealth}。
则存在 $v:Y\to\R$ 使
\[
u(t,y)\ =\ t+v(y)
\]
表示 $\succeq$。
反之，任何由 $u(t,y)=t+v(y)$ 形式表示的偏好都满足三个计价物条件。
\end{theorem}

\begin{proof}
\textbf{正方向。}
\textbf{第 1 步（定义 $v$）}：固定基准 $\bar y\in Y$。
由条件~\ref{cond:compensable}，对每个 $y\in Y$ 存在 $s\in\R$ 使
$(0,y)\sim(s,\bar y)$；此 $s$ \textbf{唯一}——
若 $(s,\bar y)\sim(0,y)\sim(s',\bar y)$，传递性给出 $(s,\bar y)\sim(s',\bar y)$，
由条件~\ref{cond:num-valuable} 得 $s=s'$。定义
\[
v(y)\ \coloneqq\ \text{唯一的}\ s\ \text{使}\ (0,y)\sim(s,\bar y).
\]

\textbf{第 2 步（把任意点折算到 $\bar y$ 行）}：
对任意 $(t,y)$，由 $(0,y)\sim(v(y),\bar y)$ 与条件~\ref{cond:no-wealth}
（取 $d=t$）得
\[
(t,\,y)\ \sim\ \bigl(t+v(y),\,\bar y\bigr).
\]

\textbf{第 3 步（表示性验证）}：
对任意 $(t,y),(t',y')\in X$，
\[
\begin{aligned}
(t,y)\succeq(t',y')
&\iff \bigl(t+v(y),\bar y\bigr)\succeq\bigl(t'+v(y'),\bar y\bigr)
   && (\text{第 2 步} + \text{传递性})\\
&\iff t+v(y)\ \ge\ t'+v(y')
   && (\text{条件~\ref{cond:num-valuable}}).
\end{aligned}
\]
即 $u(t,y)=t+v(y)$ 表示 $\succeq$。

\textbf{逆方向。}设 $u(t,y)=t+v(y)$ 表示 $\succeq$，逐条验证：
\textbf{条件~\ref{cond:num-valuable}}：$(t,y)\succeq(t',y)\iff t+v(y)\ge t'+v(y)\iff t\ge t'$。
\textbf{条件~\ref{cond:compensable}}：取 $t=v(y)-v(y')$，则
$u(0,y)=v(y)=t+v(y')=u(t,y')$，即 $(0,y)\sim(t,y')$。
\textbf{条件~\ref{cond:no-wealth}}：$(t,y)\succeq(t',y')\iff t+v(y)\ge t'+v(y')
\iff (t+d)+v(y)\ge(t'+d)+v(y')\iff(t+d,y)\succeq(t'+d,y')$。
\end{proof}

\begin{remark}[证明的结构]
三步证明做了一件事：\textbf{用条件~\ref{cond:compensable} 把每个 $y$ “定价”为 $v(y)$，
用条件~\ref{cond:no-wealth} 把定价沿计价物轴平移到任何财富水平，
用条件~\ref{cond:num-valuable} 在单一 $\bar y$ 行上完成所有比较}。
三条缺一，证明即断：无~\ref{cond:compensable} 则 $v$ 无定义；
无~\ref{cond:no-wealth} 则定价随财富漂移；
无~\ref{cond:num-valuable} 则 $\bar y$ 行上的比较无法数值化。
\end{remark}

\begin{proposition}[$v$ 的唯一性]\label{prop:ql-unique}
若 $t+v(y)$ 与 $t+\tilde v(y)$ 表示同一 $\succeq$，则 $\tilde v=v+\text{常数}$。
\end{proposition}

\begin{proof}
两者都表示时，$(0,y)\sim(0,y')\iff v(y)=v(y')\iff\tilde v(y)=\tilde v(y')$；
且 $(t,y)\sim(0,y')\iff t=v(y')-v(y)\iff t=\tilde v(y')-\tilde v(y)$，
故 $v(y')-v(y)=\tilde v(y')-\tilde v(y)$ 对一切 $y,y'$ 成立，
固定 $y$ 即得 $\tilde v-v$ 为常数。
\end{proof}

\begin{remark}[序数性辨析]
定理说“存在 $t+v(y)$ 形式的表示”，并\textbf{不}排除其他形式的表示：
$u(t,y)=t+v(y)$ 经严格递增变换 $f$ 后一般不再是拟线性形式，
却仍表示同一偏好（命题~\ref{prop:ordinal}）。
“拟线性”是偏好的\textbf{序性质}（存在某表示具此形式），
不是某个特定函数的\textbf{形式性质}。
\end{remark}

\subsection{两大后果}

\notehead{后果一：需求无财富效应（收入扩张路径平行于计价物轴）}

设 $Y=\R_+$（记其元素为 $q$），$v$ 可微、严格凹，价格结构为“每单位 $q$ 兑换 $p$ 单位计价物”，
预算约束 $t+pq=m$。效用最大化：
\[
\max_{q\ge0}\ \bigl[m-pq+v(q)\bigr]
\ \Longrightarrow\ v'(q^*)=p\quad(\text{内点解}).
\]
一阶条件\textbf{不含 $m$}：最优 $q^*$ 只由 $p$ 决定，收入 $m$ 的变化全部吸收到 $t$ 上。
收入扩张路径是一条竖直线（$q$ 恒为 $q^*$）——这就是“无财富效应”的需求侧表达。

\begin{center}
\begin{tikzpicture}[scale=1.05]
  \draw[axisline] (0,0) -- (6.2,0) node[below left, axislabel] {$q$};
  \draw[axisline] (0,0) -- (0,4.8) node[below left, axislabel] {$t$};
  % parallel ICs: t = c - sqrt(q)
  \draw[thick, MidBlue, domain=0.02:5.8, smooth, variable=\x] plot ({\x},{3.6-sqrt(\x)});
  \draw[thick, MidBlue, domain=0.02:5.8, smooth, variable=\x] plot ({\x},{2.6-sqrt(\x)});
  \draw[thick, MidBlue, domain=0.02:5.8, smooth, variable=\x] plot ({\x},{1.6-sqrt(\x)});
  \node at (5.0,1.6) {$u=c$};
  % vertical expansion path
  \draw[dashed, thick, AxisBlue] (2.6,0) -- (2.6,4.6) node[above, axislabel] {收入扩张路径};
  \node[nodeblue] at (2.6,3.6-1.612) {};
  \node[nodeblue] at (2.6,2.6-1.612) {};
  \node[nodeblue] at (2.6,1.6-1.612) {};
  \node[align=left] at (4.6,3.9) {无差异曲线互为\textbf{纵向平移}：\\ $t+v(q)=c\iff t=c-v(q)$，\\ 同一 $q$ 处斜率恒为 $v'(q)$};
\end{tikzpicture}
\end{center}
\figexplan{三条无差异曲线互为纵向平移：$t+v(q)=c$ 等价于 $t=c-v(q)$，故在同一 $q$ 处斜率恒为 $v'(q)$。竖直虚线是收入扩张路径：一阶条件 $v'(q^*)=p$ 不含收入 $m$，收入变化全部由计价物 $t$ 吸收，$q^*$ 不变——这就是无财富效应的几何表现。}

\begin{remark}[对照：正常品的收入扩张路径]
一般偏好（如 $u(q,t)=q^{1/2}t^{1/2}$）下，收入增加时 $q^*$ 与 $t^*$ 同时增加，
收入扩张路径弯曲向右上——存在财富效应。
拟线性是一般偏好的\textbf{局部近似}：当计价物支出占财富比重很大、
$q$ 的支出占比很小时，财富效应对 $q$ 的需求影响可忽略，
拟线性是“部分均衡分析”的效用基础。
\end{remark}

\notehead{后果二：福利的精确货币度量}

比较两个 $Y$ 方案 $y,y'$ 时，拟线性偏好给出\textbf{与财富无关}的补偿度量：
\[
(t,y)\succeq(t,y')\ \Longleftrightarrow\ v(y)\ge v(y'),
\]
且从 $y'$ 切换到 $y$ 的\textbf{补偿变化}恰为
\[
v(y)-v(y')\ =\ \text{使}\ (t,y')\sim\bigl(t-\bigl[v(y)-v(y')\bigr],\,y\bigr)\ \text{成立的计价物数量}.
\]
即：$v(y)-v(y')$ 是“为保持在 $y'$ 原状，愿意为换到 $y$ 支付的最大计价物数量”，
与初始财富 $t$ \textbf{无关}。一般偏好下补偿变化依赖财富水平，
$v$ 之差与财富无关正是拟线性被广泛用于福利计算（剩余度量）的原因。

\begin{example}[补偿变化的数值计算]\label{ex:ql-welfare}
设 $v(y)=\sqrt{y}$（$y$ 为某种公共品供给量），计价物 $t$ 为税收（可负，即补贴）。
比较两个方案：现状 $y'=4$（$v(y')=2$）与拟建项目 $y=9$（$v(y)=3$）。
由上述结论，从 $y'$ 切换到 $y$ 的补偿变化为
\[
v(y)-v(y')=3-2=1,
\]
即：为使纳税人在“维持现状 $y'$”与“接受项目 $y$”间无差异，需向其补偿 $1$ 单位计价物；
该数值与初始财富 $t$ 无关。若改用一般（非拟线性）偏好 $u(t,y)=t^{1/2}y^{1/2}$ 作比较，
对同一对 $(t,4)$ 与 $(t-1,9)$，无差异要求
$t^{1/2}\cdot2=(t-1)^{1/2}\cdot3$，解得 $t=9/5=1.8$ 时才无差异——
补偿量随 $t$ 变化，直观显示了财富效应如何破坏“精确货币度量”。
\end{example}

\subsection{拟线性与可分性的关系}

拟线性是可加可分（§3.6）的特例加限制：
$u(t,y)=t+v(y)$ 显然是可加形式 $u_1(t)+u_2(y)$，
额外要求 $u_1$ 为\textbf{恒等映射}（对计价物坐标“价值=数量”）。
在条件语言中：坐标独立性 + 矩形条件给出 $a\cdot u_1(t)+v(y)$ 型表示，
条件~\ref{cond:num-valuable} 与 \ref{cond:no-wealth} 进一步把 $u_1$ 钉死为线性
（无财富效应排除了 $u_1(t)=\sqrt{t}$ 之类的弯曲刻度）。
因此拟线性偏好自动满足坐标独立性与矩形条件：
$(t,y)\succeq(t',y)\iff t\ge t'$ 与 $(t,y')\succeq(t',y')$ 同真同假——
第一坐标的比较与 $y$ 固定在何处无关。

\subsection{可微情形 \texorpdfstring{$\star$}{*}}

设 $Y=\R^{n-1}_+$，$v$ 二阶连续可微。$u(t,y)=t+v(y)$ 的梯度为 $(1,\nabla v(y))$，
边际替代率
\[
\mathrm{MRS}_{t,i}(t,y)=\frac{\partial u/\partial y_i}{\partial u/\partial t}=\frac{\partial v}{\partial y_i}(y)
\]
\textbf{与 $t$ 无关}——无财富效应的可微表达。
拟凹性等价于 $v$ 凹：
$u$ 的海塞矩阵为 $\begin{pmatrix}0&0\\0&D^2v\end{pmatrix}$，
其加边海塞行列式条件归结为 $D^2v$ 负半定。
凸偏好（命题~\ref{prop:convex-qc}）$\iff$ $v$ 凹，
严格凸 $\iff$ $v$ 严格凹——拟线性把偏好凸性完全压缩为 $v$ 的凹性。
